% ══════════════════════════════════════════════════════════════════
\chapter{Introduction}
\label{ch:introduction}

The simulation of quantum systems lies at the heart of quantum computing's
promise. As Feynman observed in 1982, classical computers face a fundamental
barrier when modelling quantum mechanical systems: the Hilbert space of an
$n$-particle system grows exponentially, rendering exact classical computation
intractable for even modest system sizes~\cite{feynman1982}. Quantum computers,
by contrast, naturally encode and manipulate quantum states, offering a path
toward efficient simulation of physical systems that are otherwise
inaccessible~\cite{nielsen2010}.

This thesis applies variational quantum algorithms to the spectroscopy of
charmonium, bound states of a charm quark and its antiquark, including
spin-dependent interactions that produce the experimentally observed mass
splittings between states such as the $\eta_c$ and $J/\psi$ mesons. The
quarkonium system is physically rich enough to probe the capabilities of
near-term quantum algorithms across multiple quantum number channels and excited
states, yet small enough to be tractable with only a few qubits.

The remainder of this chapter introduces the two foundations on which the
algorithmic chapters rest: the fundamentals of quantum computing
(\cref{sec:qc}) and the variational paradigm that makes near-term quantum
hardware useful (\cref{sec:variational}). The charmonium system and its
discretisation into a finite-dimensional matrix problem are developed in
\cref{ch:framework}.


% ──────────────────────────────────────────────────────────────────
\section{Quantum Computing}
\label{sec:qc}

\subsection{Why quantum computing?}
\label{sec:why_qc}

The classical theory of computation, formalized by Turing and Church in the
1930s, rests on the assumption that information is processed according to
classical physics. For decades this abstraction posed no practical limitation.
However, the simulation of quantum mechanical systems reveals a fundamental
barrier: the state space of an $n$-particle quantum system is a Hilbert space of
dimension $2^n$, so that a faithful classical representation requires resources
that grow exponentially with system size. A modest system of 50 interacting
spin-$\tfrac{1}{2}$ particles, for instance, demands $2^{50}\approx 10^{15}$
complex amplitudes, which exceeds the memory of any existing classical computer.

In 1982, Feynman proposed a radical alternative: rather than simulating quantum
mechanics on a classical machine, one should use a controllable quantum system to
simulate another~\cite{feynman1982}. The exponential scaling that makes classical
simulation intractable becomes an \emph{asset}: a quantum processor with $n$
qubits naturally inhabits the same $2^n$-dimensional space as the target system.
This observation laid the conceptual foundation for quantum computing.

The subsequent development of quantum algorithms by Deutsch, Shor, and
Grover~\cite{nielsen2010} demonstrated that quantum computers are not merely
faster classical machines, but represent a genuinely different model of
computation. The complexity class $\mathsf{BQP}$ (bounded-error quantum
polynomial time) is believed to be strictly larger than its classical counterpart
$\mathsf{BPP}$, meaning there probably exist problems for which quantum algorithms
achieve an exponential advantage that no classical algorithm can
match~\cite{nielsen2010}. We highlight the former has not been proved yet, 
despite a long list of attempts by mathematicians. The source of this advantage lies in two uniquely
quantum phenomena: \emph{superposition}, which allows exploration of exponentially
many computational paths simultaneously, and \emph{interference}, which enables
the selective amplification of correct answers and suppression of incorrect ones.


\subsection{Qubits and quantum states}
\label{sec:qubits}

The fundamental unit of quantum information is the \emph{qubit}, a two-level
quantum system whose pure state is described by a normalized vector in
$\mathbb{C}^2$:
\begin{equation}
  \label{eq:qubit}
  \ket{\psi} = \alpha\ket{0} + \beta\ket{1},
  \qquad |\alpha|^2 + |\beta|^2 = 1,
\end{equation}
where $\alpha,\beta\in\mathbb{C}$ are \emph{probability amplitudes} and
$\{\ket{0},\ket{1}\}$ is the computational basis (conventionally identified with
the eigenstates of the Pauli $Z$ operator). Unlike a classical bit, which is
deterministically~0 or~1, a qubit exists in a coherent \emph{superposition} of
both basis states.

Geometrically, any single-qubit pure state can be represented as a point on the
\emph{Bloch sphere} via the parameterization
\begin{equation}
  \label{eq:bloch}
  \ket{\psi}
    = \cos\frac{\theta}{2}\,\ket{0}
    + e^{i\phi}\sin\frac{\theta}{2}\,\ket{1},
\end{equation}
with polar angle $\theta\in[0,\pi]$ and azimuthal angle $\phi\in[0,2\pi)$. The
north and south poles correspond to $\ket{0}$ and $\ket{1}$, respectively, while
points on the equator represent equal superpositions with varying relative phase.

For a system of $n$ qubits, the state space is the tensor product
$\mathcal{H} = (\mathbb{C}^2)^{\otimes n}$, spanned by $2^n$ computational
basis states $\ket{x_1 x_2 \cdots x_n}$ with $x_i\in\{0,1\}$. A general pure
state is therefore
\begin{equation}
  \label{eq:nqubit}
  \ket{\Psi} = \sum_{x\in\{0,1\}^n} c_x \ket{x},
  \qquad \sum_x |c_x|^2 = 1,
\end{equation}
requiring $2^n$ complex amplitudes to specify. It is this exponential growth of
the state space that makes quantum systems both hard to simulate classically and
powerful as computational devices.


\subsection{Entanglement}
\label{sec:entanglement}

A multi-qubit state is called \emph{separable} (or a product state) if it can be
written as a tensor product of individual qubit states,
$\ket{\Psi} = \ket{\psi_1}\otimes\ket{\psi_2}\otimes\cdots\otimes\ket{\psi_n}$.
States that \emph{cannot} be decomposed in this way are \emph{entangled}. The
simplest example is the Bell state
\begin{equation}
  \label{eq:bell}
  \ket{\Phi^+} = \frac{1}{\sqrt{2}}\bigl(\ket{00}+\ket{11}\bigr),
\end{equation}
which describes two qubits in a joint superposition: neither qubit has a definite
value, yet their measurement outcomes are perfectly correlated. This correlation
persists regardless of the spatial separation between the qubits and cannot be
explained by any local hidden variable theory, as demonstrated by Bell's
theorem~\cite{bell1964}.

Entanglement is one of the biggest computational resources that quantum computers give us.
Classical correlations can be shared between bits using shared
randomness, but entangled qubits carry correlations that are strictly stronger.
In quantum algorithms, entangling gates create superpositions over correlated
multi-qubit configurations, enabling the exploration of an exponentially large
solution space in a way that has no classical counterpart. Without entanglement, a
quantum computer can be efficiently simulated classically~\cite{nielsen2010}.


\subsection{Quantum gates and circuits}
\label{sec:gates}

The time evolution of a closed quantum system is governed by a unitary operator.
Quantum computation exploits this by decomposing the desired evolution into a
sequence of elementary \emph{quantum gates}: unitary transformations acting on one
or two qubits at a time.

The three Pauli operators form a basis for single-qubit operations:
\begin{equation}
  \label{eq:pauli}
  X = \begin{pmatrix} 0 & 1 \\ 1 & 0 \end{pmatrix},\quad
  Y = \begin{pmatrix} 0 & -i \\ i & 0 \end{pmatrix},\quad
  Z = \begin{pmatrix} 1 & 0 \\ 0 & -1 \end{pmatrix}.
\end{equation}
Together with the identity $I$, they satisfy $\sigma_i^2 = I$, anticommute
pairwise ($\sigma_i\sigma_j = -\sigma_j\sigma_i$ for $i\neq j$), and form a
basis for the space of $2\times 2$ Hermitian matrices. Any $2\times 2$ Hermitian
operator, and hence any single-qubit observable, can be written as
$A = a_0 I + a_1 X + a_2 Y + a_3 Z$ with $a_i\in\mathbb{R}$.

The Pauli operators generate continuous families of rotation gates on the Bloch
sphere:
\begin{equation}
  \label{eq:rotations}
  R_\mu(\theta) = e^{-i\theta\,\mu/2}
    = \cos\frac{\theta}{2}\,I - i\sin\frac{\theta}{2}\,\mu,
  \qquad \mu \in \{X,Y,Z\}.
\end{equation}
The Hadamard gate $H = (X+Z)/\sqrt{2}$ maps $\ket{0}\mapsto(\ket{0}+\ket{1})/\sqrt{2}$,
creating an equal superposition from a computational basis state.

Two-qubit entangling gates introduce correlations between qubits. The
controlled-NOT (CNOT) gate flips the target qubit conditioned on the control
qubit being in state $\ket{1}$:
\begin{equation}
  \label{eq:cnot}
  \text{CNOT} = \ket{0}\bra{0}\otimes I + \ket{1}\bra{1}\otimes X
  = \begin{pmatrix}
    1 & 0 & 0 & 0 \\
    0 & 1 & 0 & 0 \\
    0 & 0 & 0 & 1 \\
    0 & 0 & 1 & 0
  \end{pmatrix}.
\end{equation}
The set $\{R_X,R_Y,R_Z,\text{CNOT}\}$ constitutes a \emph{universal gate set}:
any unitary on $n$ qubits can be approximated to arbitrary precision by a finite
sequence of these gates~\cite{nielsen2010}.

A quantum algorithm is specified as a \emph{quantum circuit}: a sequence of gates
applied to a register of qubits initialized in a known state, followed by
measurement. \Cref{fig:bell_circuit} illustrates a minimal circuit that prepares
the Bell state of \cref{eq:bell} using only a Hadamard gate and a CNOT.

\begin{figure}[htbp]
\centering
\begin{quantikz}
  \lstick{$\ket{0}$} & \gate{H} & \ctrl{1} & \meter{} \\
  \lstick{$\ket{0}$} & \qw      & \targ{}  & \meter{}
\end{quantikz}
\caption{Quantum circuit preparing the Bell state
  $\ket{\Phi^+}=(\ket{00}+\ket{11})/\sqrt{2}$. The Hadamard gate creates a
  superposition on the first qubit; the CNOT entangles both qubits. Measurement
  yields $\ket{00}$ or $\ket{11}$, each with probability $1/2$.}
\label{fig:bell_circuit}
\end{figure}


\subsection{Quantum parallelism and interference}
\label{sec:parallelism}

The mechanism behind quantum computational advantage is often described loosely
as ``computing on all inputs at once.'' While evocative, this picture is
incomplete: the key ingredient is not parallelism alone, but the ability to make
the parallel branches \emph{interfere}.

Consider a unitary $U_f$ that evaluates a function $f:\{0,1\}^n\to\{0,1\}$.
Applied to an equal superposition over all $2^n$ inputs, it produces
\begin{equation}
  \label{eq:parallelism}
  U_f \frac{1}{\sqrt{2^n}}\sum_{x=0}^{2^n-1}\ket{x}\ket{0}
  = \frac{1}{\sqrt{2^n}}\sum_{x=0}^{2^n-1}\ket{x}\ket{f(x)}.
\end{equation}
All function values are encoded in a single quantum state. However, measuring
this state yields a \emph{single} random pair $(x,f(x))$ - no better than
classical evaluation. The power of quantum algorithms lies in subsequent
operations that exploit \emph{interference}: unitary gates are designed so that
amplitudes corresponding to desired outputs add constructively, while those
corresponding to undesired outputs cancel destructively. The final measurement
then yields the correct answer with high probability.

This interplay of superposition, entanglement, and interference is the
mathematical core of quantum speedups. Superposition provides exponential
parallelism; entanglement enables correlations across qubits that have no
classical equivalent; and interference channels the computation toward the
correct outcome.


\subsection{Measurement and the Born rule}
\label{sec:measurement}

While unitary evolution is deterministic, extracting information from a quantum
system is inherently probabilistic. A \emph{projective measurement} in the
computational basis is described by the set of projectors
$\{\Pi_x = \ket{x}\bra{x}\}_{x\in\{0,1\}^n}$. Given a state
$\ket{\Psi}=\sum_x c_x\ket{x}$, the \emph{Born rule} prescribes the outcome
probabilities:
\begin{equation}
  \label{eq:born}
  p(x) = |\braket{x}{\Psi}|^2 = |c_x|^2.
\end{equation}
Upon obtaining outcome $x$, the state collapses to $\ket{x}$, and any
information encoded in the remaining amplitudes is irretrievably lost.

This has a direct practical consequence: a single measurement of $n$ qubits
yields only $n$ classical bits, despite the state containing $2^n$ amplitudes.
To estimate a continuous quantity, such as an energy expectation value, the
measurement must be repeated many times ($N_{\text{shots}}$ repetitions), and the
result is subject to statistical uncertainty scaling as $1/\sqrt{N_{\text{shots}}}$.

More generally, to measure an observable $O$ (a Hermitian operator), its
expectation value in state $\ket{\Psi}$ is given by
\begin{equation}
  \label{eq:expectation}
  \expval{O} = \bra{\Psi}O\ket{\Psi} = \sum_i \lambda_i\, p_i,
\end{equation}
where $\lambda_i$ are the eigenvalues of $O$ and $p_i$ the probability of
obtaining the corresponding eigenstate. The concrete measurement protocol
for a Pauli-decomposed Hamiltonian is constructed in \cref{ch:framework}.


\subsection{Density matrices and mixed states}
\label{sec:density}

The pure state formalism of \cref{eq:qubit} describes isolated, perfectly
controlled quantum systems. Real quantum processors, however, interact with their
environment, leading to \emph{decoherence} and \emph{mixed states} that require a
more general description.

The \emph{density matrix} (or density operator) of a pure state
$\ket{\psi}$ is
\begin{equation}
  \label{eq:density_pure}
  \rho = \ket{\psi}\bra{\psi}.
\end{equation}
A \emph{mixed state} is described by a convex combination
\begin{equation}
  \label{eq:density_mixed}
  \rho = \sum_i p_i \ket{\psi_i}\bra{\psi_i},
  \qquad p_i \geq 0,\quad \sum_i p_i = 1.
\end{equation}
The density matrix is Hermitian, positive semidefinite, and has unit trace. A
state is pure if and only if $\mathrm{Tr}(\rho^2) = 1$; for mixed states,
$\mathrm{Tr}(\rho^2) < 1$, quantifying the degree of decoherence.

In the density matrix formalism, unitary evolution becomes
$\rho \mapsto U\rho\, U^\dagger$, and expectation values are computed as
$\expval{O}=\mathrm{Tr}(O\rho)$. Noise process, such as the depolarizing
channel relevant to this thesis, are described by completely positive
trace-preserving (CPTP) maps on $\rho$. For example, the single-qubit
depolarizing channel with error rate $p$ acts as
\begin{equation}
  \label{eq:depolarizing}
  \mathcal{E}(\rho) = (1-p)\,\rho + \frac{p}{3}
    \bigl(X\rho X + Y\rho Y + Z\rho Z\bigr),
\end{equation}
which with probability $p$ replaces the state with a random Pauli error.
The depolarising channel is the workhorse of the noise studies carried out
later in the thesis: it is analytically tractable and captures the dominant
isotropic effect of incoherent gate errors on present-day quantum hardware.


\subsection{Quantum simulation}
\label{sec:qsim}

Among the many proposed applications of quantum computing, including
chemestry, optimization, and machine learning, the simulation of quantum
systems stands out as the most natural and arguably the most promising near-term
application~\cite{feynman1982,preskill2018}. The reason is both conceptual and
practical: a quantum computer \emph{is} a quantum system, so using it to simulate
another quantum system leverages its native capabilities rather than working
against them.

The essential mapping is straightforward. A physical system with a
$2^n$-dimensional Hilbert space is encoded on $n$ qubits. The system's
Hamiltonian $H$ is decomposed into a sum of Pauli operators acting on these
qubits:
\begin{equation}
  \label{eq:pauli_ham}
  H = \sum_{i=1}^{M} c_i\, P_i,
  \qquad P_i \in \{I,X,Y,Z\}^{\otimes n},
\end{equation}
where each $P_i$ is a tensor product of single-qubit Pauli matrices (a
\emph{Pauli string}) and $c_i\in\mathbb{R}$ are real coefficients. This
decomposition is always possible since the $4^n$ Pauli strings form a complete
basis for $2^n\times 2^n$ Hermitian matrices.

Ground state energies, excited state spectra, and dynamical properties can then
be extracted using quantum algorithms that prepare and measure appropriate states
on the quantum register. For the small systems considered in this thesis, the
encoding requires only 2--3 qubits, but the same framework scales to larger
systems where classical methods become intractable.

Quantum simulation of atomic, molecular, and nuclear systems has seen rapid
progress in recent years, with demonstrations ranging from small molecules to
lattice gauge theories~\cite{peruzzo2014}. The quarkonium system studied in this
thesis belongs to this broader programme, representing a step toward quantum
simulation of QCD bound states.


\subsection{The NISQ era}
\label{sec:nisq}

Current quantum hardware operates in the \emph{Noisy Intermediate-Scale Quantum}
(NISQ) regime~\cite{preskill2018}, characterized by processors with tens to
hundreds of qubits subject to significant imperfections:
\begin{itemize}
  \item \textbf{Gate errors} on the order of $10^{-3}$--$10^{-2}$ per operation,
    described by noise channels such as the depolarizing channel of
    \cref{eq:depolarizing};
  \item \textbf{Limited connectivity} between qubits, requiring additional SWAP
    gates that increase circuit depth;
  \item \textbf{Short coherence times} that bound the number of operations
    executable before decoherence degrades the quantum state.
\end{itemize}
These constraints preclude the deep circuits required by fault-tolerant algorithms
such as quantum phase estimation, which need thousands of controlled operations
and ancilla qubits for error correction.

The NISQ era has therefore driven the development of algorithms specifically
designed for shallow circuits and resilient to moderate noise levels. The most
successful paradigm to date is the \emph{variational} approach, in which a
quantum processor prepares parameterized trial states while a classical computer
optimizes the parameters.


% ──────────────────────────────────────────────────────────────────
\section{Variational Quantum Methods}
\label{sec:variational}

The NISQ constraints of \cref{sec:nisq} rule out algorithms that demand
deep, fully coherent circuits. Variational quantum algorithms sidestep
this restriction by offloading the search itself onto a classical
optimiser: the quantum processor is asked only to prepare a parameterised
state $\ket{\psi(\bm{\theta})} = U(\bm{\theta})\ket{0}^{\otimes n}$ and to
estimate the expectation value
$E(\bm{\theta}) = \bra{\psi(\bm{\theta})}H\ket{\psi(\bm{\theta})}$, while
a classical routine drives $\bm{\theta}$ toward the minimum. The loop is
shown in \cref{fig:vqe_loop}; its formal derivation from the
Rayleigh--Ritz principle and the explicit measurement protocol on Pauli
strings are developed in \cref{ch:framework}, where they specialise to
each of the three algorithms studied in this thesis.

\begin{figure}[htbp]
\centering
\begin{tikzpicture}[
    node distance=1.6cm and 2.5cm,
    box/.style={
      draw, rounded corners=4pt, minimum width=3.2cm,
      minimum height=1cm, align=center, font=\small
    },
    arrow/.style={-{Stealth[length=6pt]}, thick},
  ]
  \node[box, fill=blue!8] (prep) {Prepare\\$\ket{\psi(\bm{\theta})}$};
  \node[box, fill=blue!8, right=of prep] (meas)
    {Measure\\$\expval{P_i}_{\bm{\theta}}$};
  \node[box, fill=orange!10, below=of meas] (eval)
    {Compute\\$E(\bm{\theta})=\sum c_i\expval{P_i}$};
  \node[box, fill=orange!10, below=of prep] (opt)
    {Classical optimiser\\update $\bm{\theta}$};

  \draw[arrow] (prep) -- (meas);
  \draw[arrow] (meas) -- (eval);
  \draw[arrow] (eval) -- (opt);
  \draw[arrow] (opt) -- (prep);

  \node[font=\footnotesize\itshape, above=0.2cm of prep, anchor=south west,
        xshift=-0.3cm] {Quantum processor};
  \node[font=\footnotesize\itshape, below=0.2cm of opt, anchor=north west,
        xshift=-0.3cm] {Classical computer};

  \draw[dashed, gray] ($(prep.south west)+(-0.7,-0.8)$) --
                       ($(meas.south east)+(0.7,-0.8)$);
\end{tikzpicture}
\caption{The hybrid variational loop. The quantum processor prepares a
  parameterised state and estimates the Pauli expectation values that
  compose the energy; the classical optimiser proposes the next
  $\bm{\theta}$. The cost function is a statistical average over circuit
  shots, which is what makes the loop robust to the gate errors of NISQ
  hardware.}
\label{fig:vqe_loop}
\end{figure}

The remainder of this section explains \emph{why} this hybrid arrangement
is the right framework for present-day hardware, and where its limits
lie. Two properties are responsible for both its success and its
characteristic failure modes: the noise resilience of an averaged cost
function, and the tension between the expressibility of a parameterised
ansatz and its trainability.

\subsection{Why the variational principle survives noise}
\label{sec:vqa_noise}

The decisive feature of a variational algorithm is that the quantity it
optimises is a statistical average. A single circuit shot returns a
bit-string. The cost $E(\bm{\theta})$ is the empirical mean of many such
shots, weighted by the Pauli coefficients of $H$. Stochastic shot noise
and biased gate errors therefore enter the optimiser as a deformation of
the cost landscape, not as a corrupted coherent observable that must
survive intact to the final measurement.

Two consequences follow. First, depolarising-like noise of strength $p$
applied uniformly across the circuit damps every nontrivial Pauli
expectation by a common factor that depends on the circuit structure but
not on $\bm{\theta}$ once the optimiser sits near a smooth minimum. The
position of the minimum in parameter space is preserved, while the
absolute energy at that minimum is biased upward by an amount that can be
estimated and partially removed by zero-noise extrapolation. Second, the
optimiser is insensitive to constant offsets in the cost: any
$\bm{\theta}$-independent shift cancels under finite differences and
gradient updates. This is the formal sense in which variational
algorithms are self-correcting against coherent miscalibration, and it is
the property that justifies their use on devices whose two-qubit gate
infidelities sit orders of magnitude above the fault-tolerance
threshold~\cite{mcclean2016,preskill2018}.

The contrast with quantum phase estimation is informative. Phase
estimation reads the ground-state energy off the relative phase of an
ancilla, accumulated coherently over a sequence of controlled-$U$
applications whose length scales with the desired precision. A single
decohering event in that sequence corrupts the phase and the answer.
Variational methods replace this brittle coherent encoding by an
incoherent average, paying a $1/\sqrt{N_{\mathrm{shots}}}$ statistical
cost in exchange for graceful degradation under the noise budget of NISQ
hardware.

\subsection{Expressibility and trainability}
\label{sec:vqa_expressibility}

The flip side of this resilience is that the answer is only as good as
the ansatz. Two notions sharpen the question of how to choose
$U(\bm{\theta})$.

\emph{Expressibility} measures the extent to which the parameterised
manifold $\{U(\bm{\theta})\ket{0}^{\otimes n} : \bm{\theta}\in\mathbb{R}^p\}$
covers the physically relevant subspace of Hilbert space. If the true
ground state lies outside this manifold, the variational minimum is a
strict upper bound on $E_0$, and that bound cannot be tightened by more
optimiser iterations or more shots. Hardware-efficient ansatze chase
expressibility through depth: every additional layer of entangling gates
and single-qubit rotations enlarges the reachable manifold.

\emph{Trainability} measures the conditioning of the cost landscape.
A maximally expressive ansatz on $n$ qubits suffers the barren-plateau
phenomenon~\cite{mcclean2018,cerezo2021}: the variance of the energy
gradient with respect to any single parameter decays exponentially in $n$,
so the landscape is essentially flat across the bulk of parameter space
and the optimiser receives no signal. Pushing for more expressibility
through brute depth therefore destroys trainability, yielding ansatze
that can in principle reach the ground state but cannot be steered there.

The way out is to give up universal coverage in favour of
problem-specific symmetries. A particle-number-conserving manifold, for
instance, restricts the search to the physically meaningful sector and
shrinks the parameter space by orders of magnitude, restoring nonzero
gradients. The Unitary Coupled Cluster ansatz used by Gallimore and Liao
(\cref{sec:gallimore}) and the three-parameter $R_y$ ansatz used by
Woloshyn (\cref{sec:woloshyn}) are both instances of this strategy: each
is tailored to its qubit encoding so that the trial manifold contains
the ground state of the channel Hamiltonian while remaining
low-dimensional enough to be trainable on current hardware.

Beyond the ground state, the same hybrid framework extends naturally:
orthogonality penalties give access to excited states, imaginary-time
equations of motion give a noise-stable route to the ground state from
arbitrary starting points, and parameterised propagators give dynamical
observables. The three algorithms developed in the next chapter cover
precisely these three extensions on the same charmonium Hamiltonian,
making the strengths and limits of the variational framework visible in
a single physical system.


% ──────────────────────────────────────────────────────────────────
\section{Thesis Overview}
\label{sec:overview}

This thesis applies three variational quantum algorithms to charmonium
spectroscopy, comparing them on the same ${}^1S_0$, ${}^3S_1$, and
${}^1P_1$ channel Hamiltonians and on the same noise model. The goal is
to test whether the variational framework introduced in this chapter can
recover spin-dependent mass splittings of the charmonium spectrum on
hardware whose fidelity is firmly within the NISQ regime, and to
identify which features of each algorithm carry that performance.

The remaining chapters are organised as follows:
\begin{itemize}
  \item \textbf{Chapter~2} introduces the charmonium system, derives the
    radial Hamiltonian in a truncated harmonic-oscillator basis, and
    develops in self-contained form the three algorithms studied in this
    thesis: the Variational Quantum Eigensolver of Gallimore and Liao,
    the Variational Imaginary-Time Evolution method of Woloshyn, and the
    Trotter--Taylor Imaginary-Time Evolution scheme of Yi et al.
  \item \textbf{Chapter~3} presents the numerical results: noiseless
    benchmarks of energies and transition amplitudes, and the
    noisy-simulator study that quantifies how each algorithm degrades
    under depolarising gate noise.
\end{itemize}

